% ============================================================
\section{ASR}\label{sec:asr}
% ============================================================

Un requisito arquitectónicamente significativo (ASR) es un requisito cuya satisfacción tiene un efecto profundo sobre la arquitectura del sistema, de manera que la arquitectura podría ser diferente si ese requisito no existiera. Esta sección identifica los ASR de Bastion, dice de dónde procede cada uno y los prioriza. Se construye solo a partir de los documentos de entrada.

% ------------------------------------------------------------
\subsection{Cómo se identifican}

\textbf{La prueba.} A cada requisito candidato se le hace la misma pregunta: \emph{¿la ausencia de este requisito podría llevarnos a diseñar una arquitectura significativamente diferente?} Si la respuesta es sí, se nombran las decisiones que obliga a tomar. Por ejemplo, recuperar una partida después de una falla del servidor en un máximo de 60 segundos es ASR porque obliga a decidir la persistencia del estado de la partida, cómo detectar el fallo, los mecanismos de recuperación, el despliegue y la sincronización. Si la respuesta es no, el requisito sigue siendo obligatorio, pero se cumple dentro de cualquier arquitectura razonable y no sirve para elegir entre ellas.

\textbf{Las fuentes.} La prueba se aplica a los candidatos de las tres fuentes de las que puede surgir un ASR, sin olvidar los objetivos de negocio, los stakeholders, los concerns y las dependencias externas que las respaldan.

\begin{tabla}{|L{4.0cm}|Y|}
\hline
\filacab \cab{Fuente} & \cab{Qué se revisó} \\ \hline
\endhead
Funcionalidad arquitectónicamente significativa &
Los 49 casos de uso del documento de casos de uso de Bastion. Solo cuentan los que, para realizarse, obligan a una decisión estructural: la mayoría se resuelve sobre la estructura que exigen otros, como se explica en la subsección~\ref{sec:asr:no}. \\ \hline
Escenarios de calidad &
Los 18 escenarios del documento de contexto, priorizados en su Utility Tree por importancia y sensibilidad arquitectónica. \\ \hline
Restricciones &
Las restricciones técnicas y de proyecto (CON) y los requisitos impuestos (REQ) que eliminan alternativas. Incluye la organización del propio equipo: dos personas y 14 semanas. \\ \hline
Objetivos, stakeholders, concerns y dependencias &
Los objetivos de negocio (juego en línea, público de 8 años en adelante, distribución internacional y originalidad frente a Quoridor), los 19 stakeholders con sus 26 concerns y las dependencias externas DEP-01 a DEP-04 y DEP-08. No son una fuente aparte: se usan para confirmar que el requisito es real y para medir su importancia. \\ \hline
\end{tabla}

% ------------------------------------------------------------
\subsection{ASR encontrados}

Al aplicar la prueba a las tres fuentes se encontraron quince ASR: catorce activos y uno reservado. La tabla los reúne todos, con lo que la arquitectura debe cumplir, su procedencia y los elementos del análisis de los que vienen. El orden de la tabla es el de su numeración; el orden de importancia se establece en la subsección~\ref{sec:asr:prioridad}.

\begin{tabla}{|L{1.3cm}|Y|L{2.4cm}|L{4.3cm}|}
\hline
\filacab \cab{ASR} & \cab{Qué debe cumplir la arquitectura} & \cab{Procedencia} & \cab{Viene de} \\ \hline
\endhead
ASR-01 & Una jugada en línea se responde en menos de 1 segundo, con mediana bajo 300 ms, con el servidor decidiendo y el tráfico cifrado. & Escenario, Restricción & QAS-01; REQ-02; CRN-02; STK-01, STK-09; DEP-04. \\ \hline
ASR-02 & Un cliente alterado no puede cambiar el estado de una partida ni obtener datos de otro jugador. & Escenario & QAS-05; CRN-10, CRN-23; STK-15, STK-10, STK-05. \\ \hline
ASR-03 & El tráfico capturado no revela credenciales, códigos, sesiones ni chat, y no se puede reutilizar. & Escenario, Restricción & QAS-06; REQ-03, CON-02, CON-03; CRN-22, CRN-05; STK-15, STK-02; DEP-02. \\ \hline
ASR-04 & Una desconexión no ordenada no deja la partida indefinida: se detecta en menos de 15 segundos y el jugador retoma en menos de 5 si vuelve antes de 2 minutos, incluso desde otro dispositivo. & Escenario, Funcionalidad & QAS-08; CU-26, CU-24; CON-17; CRN-06; STK-01, STK-14; DEP-04. \\ \hline
ASR-05 & Una caída de SQL Server pierde como máximo 5 segundos de jugadas y ningún dato confirmado, y deja constancia de lo perdido. & Escenario, Restricción & QAS-09; CON-04; CU-25; CRN-19; STK-13, STK-14; DEP-03, DEP-04. \\ \hline
ASR-06 & Las reglas cambian sin tocar código y una partida completa se reproduce sin interfaz. & Escenario & QAS-12, QAS-17; CON-15; CRN-13, CRN-17; STK-06, STK-10. \\ \hline
ASR-07 & Reservado: partidas en red local con el modo en línea ya entregado. & Escenario, condicionado a Q-02 & QAS-14; Q-02; CRN-20; STK-14. \\ \hline
ASR-08 & Dos o más jugadores juegan partidas completas en línea, creadas por emparejamiento o por sala, con turnos y relojes arbitrados por un servidor. & Funcionalidad, Restricción & REQ-01; CU-17, CU-18, CU-19, CU-20, CU-21, CU-25; objetivo de juego en línea; STK-03, STK-01. \\ \hline
ASR-09 & Una partida en curso se puede observar, incluso sin cuenta, recibiendo el estado con retraso y sin frenar a los jugadores. & Funcionalidad & CU-34, CU-12 (permite\_espectadores), CU-21 (paso 20); CRN-02; STK-01. \\ \hline
ASR-10 & Un jugador se enfrenta a un rival controlado por el servidor, con niveles y sugerencias, dentro del mismo presupuesto por jugada. & Funcionalidad & CU-27; Q-01; CRN-03, CRN-14; STK-06, STK-04. \\ \hline
ASR-11 & Todo mensaje de chat pasa por un filtro en el servidor que falla cerrado, y cada reporte y sanción deja un rastro revisable sin guardar conversaciones no reportadas. & Funcionalidad, Restricción, Escenario & CU-28, CU-33, CU-42 a CU-45; REQ-04, REQ-05; QAS-04, QAS-07; CRN-04, CRN-05, CRN-18; STK-02, STK-17, STK-19. \\ \hline
ASR-12 & El acceso exige segundo factor y un niño de 8 años lo completa solo o con su tutor. & Restricción, Escenario & REQ-03, REQ-07; QAS-03; CU-01, CU-04; CRN-01, CRN-26; TEN-03; STK-15, STK-02; DEP-01, DEP-02. \\ \hline
ASR-13 & Todo texto y formato se adapta a otro idioma sin tocar código, y el texto viaja en Unicode desde la pantalla hasta la base de datos. & Restricción, Escenario & REQ-06, CON-05, CON-04; QAS-11; CU-12; objetivo de distribución internacional; CRN-15, CRN-16; STK-11. \\ \hline
ASR-14 & La solución usa solo C\# sobre .NET, SQL Server como almacenamiento y TCP sin WebSockets. & Restricción & CON-01, CON-02, CON-03, CON-04, CON-16, CON-17; tecnología asignada por el cliente (STK-03). \\ \hline
ASR-15 & La solución se construye en 14 semanas, con un incremento auditable cada semana, por un equipo de dos personas, y lo que no es esencial se puede recortar sin tocar lo esencial. & Restricción, Escenario & CON-13, CON-14; QAS-13, QAS-17; equipo de dos personas; CRN-07, CRN-08, CRN-14; STK-03, STK-04, STK-16. \\ \hline
\end{tabla}

\textbf{Qué muestra la procedencia.} Once ASR se apoyan en al menos un escenario, nueve en una restricción y cinco en una funcionalidad, y nueve combinan más de una fuente.

% ------------------------------------------------------------
\subsection{Por qué cada uno es ASR}

Para cada ASR de la tabla anterior se responde la pregunta de la prueba: qué decisiones obliga a tomar y cómo podría ser la arquitectura si no existiera.

\begin{tabla}{|L{1.3cm}|Y|Y|}
\hline
\filacab \cab{ASR} & \cab{¿Es ASR? Sí, porque obliga a decidir\ldots} & \cab{Sin este requisito, la arquitectura podría\ldots} \\ \hline
\endhead
ASR-01 &
Cuántos viajes de red hay por jugada, cuánto trabajo cabe dentro del turno, qué sale del camino de la respuesta, cómo se reparte el segundo entre red, cifrado, arbitraje y notificación, y cómo se mide cada tramo. &
Validar cada jugada contra la base de datos, escribir de forma síncrona y notificar en secuencia: un diseño más simple y más lento. \\ \hline
ASR-02 &
Quién tiene la autoridad sobre el estado, dónde se ejecuta el motor de reglas, cómo se valida cada mensaje antes de tocar el estado y cómo se registra cada intento. &
Dejar que el cliente calcule y difunda el estado con el servidor como simple retransmisor, o jugar entre pares sin servidor de partidas. \\ \hline
ASR-03 &
Dónde vive el cifrado, sobre la conexión o mensaje por mensaje; cómo se gestionan certificados y claves; cómo es el apretón de manos; y cómo se impide reutilizar un token capturado. &
Hacer viajar el protocolo propio en claro, sin apretón de manos ni vencimiento de sesiones. \\ \hline
ASR-04 &
Que la sesión de partida sea distinta de la conexión TCP, cómo se detecta una conexión medio abierta, cómo se reasocia una sesión a otra conexión o a otro dispositivo, dónde vive el estado mientras el jugador no está y qué reloj mide los plazos. &
Ligar la partida al socket y darla por terminada cuando la conexión se cierra. \\ \hline
ASR-05 &
Qué se escribe de forma confirmada y qué de forma diferida, cuánto se puede perder, cómo se detecta la pérdida de la conexión con la base de datos, qué ocurre con las partidas mientras no está y cómo se reconstruyen. &
Escribir todo de forma síncrona o todo al final de la partida, sin cola, sin registro de lo perdido y sin suspensión. \\ \hline
ASR-06 &
Dónde viven los parámetros de las reglas, que el motor sea una pieza independiente de la interfaz, de la red y de la base de datos, y que sea determinista. &
Tener las reglas como constantes dentro del código del servidor, con el motor mezclado con el arbitraje y con la interfaz. \\ \hline
ASR-08 &
Que exista una arquitectura cliente-servidor, con un servicio de partidas, emparejamiento y salas, gestión de turnos y relojes, y persistencia central de cuentas y resultados. &
Ser un juego local en un solo equipo, sin red, sin servidor, sin protocolo y sin cuentas centralizadas. Es la diferencia más grande de toda la lista. \\ \hline
ASR-09 &
Cómo se difunde el estado a destinatarios que no juegan, cómo se aplica el retraso, cómo se aceptan conexiones sin sesión autenticada sin abrir una vía de ataque y cómo se evita que el número de espectadores frene a los jugadores. &
Notificar a cada partida solo a sus participantes, todos autenticados, sin ningún mecanismo de difusión ni de retraso. \\ \hline
ASR-10 &
Dónde se ejecuta la IA, cómo ocupa el lugar de un participante, cuánto de su cálculo cabe en el presupuesto por jugada y cómo se retira si el alcance se recorta. &
Suponer siempre dos personas conectadas, con un servidor que nunca calcula jugadas por su cuenta. \\ \hline
ASR-11 &
Que el filtro viva en el servidor; que el chat sea un servicio aparte con canales que exceden la partida, global y amigos; qué hace cuando el filtro no responde; qué se conserva de los mensajes; y cómo se liga cada sanción a su evidencia. &
Prescindir del servicio de chat y del de moderación, del catálogo de palabras prohibidas y de cualquier dato de menores que proteger. \\ \hline
ASR-12 &
Integrar un proveedor externo del segundo factor (DEP-02) y posiblemente el correo (DEP-01), gestionar secretos, dar a las cuentas estados propios, como pendiente de verificación, bloqueo temporal o aprobación del tutor, y aislar el mecanismo detrás de una interfaz mientras Q-03 siga abierta. &
Resolver el acceso con usuario y contraseña, con un módulo de cuentas que no depende de ningún tercero. \\ \hline
ASR-13 &
Convenciones que atraviesan todos los componentes: textos fuera del código, Unicode en la comunicación y en las columnas, formatos por región y una interfaz sin anchos fijos. No cambia qué componentes existen, pero si se descubre tarde obliga a tocar cada pantalla, cada mensaje y cada columna. &
Escribir los textos en el código, usar columnas sin Unicode y maquetar las pantallas para un solo idioma. \\ \hline
ASR-14 &
Construir un protocolo propio sobre TCP, porque con WebSockets se heredarían la delimitación, el apretón de manos y el cifrado; diseñar la persistencia sobre un modelo relacional; y compartir código entre cliente y servidor en un mismo ecosistema. &
Elegir libremente un protocolo estándar con esas capacidades ya resueltas, otro modelo de datos o un motor de juego en otro lenguaje. El protocolo dejaría de ser un componente a diseñar. \\ \hline
ASR-15 &
Construir por rebanadas verticales, probar sin interfaz, aislar solo lo que puede recortarse y mantener pocas piezas desplegables. Es la estructura de asignación de trabajo de Bass, Clements y Kazman: el tamaño del equipo condiciona en cuántas piezas conviene partir el sistema. &
Plantearse servidores replicados, alta disponibilidad de la base de datos o muchos servicios desplegados por separado, que con más tiempo o más gente serían alternativas razonables. \\ \hline
\end{tabla}

ASR-07 no aparece en la tabla porque su prueba depende de Q-02: si se aprueba la LAN, obligará a separar el servidor de partidas y su configuración del cliente; si no, no existe.

% ------------------------------------------------------------
\subsection{Requisitos y decisiones que no son ASR}\label{sec:asr:no}

La prueba también sirve para descartar. Los siguientes elementos son obligatorios o importantes, pero su ausencia no llevaría a una arquitectura significativamente diferente: se cumplen sobre la estructura que los ASR ya imponen.

\begin{tabla}{|L{4.2cm}|L{2.3cm}|Y|}
\hline
\filacab \cab{Elemento} & \cab{Fuente} & \cab{Por qué no pasa la prueba} \\ \hline
\endhead
Perfil, nickname y cosméticos (CU-12 a CU-16) & Funcionalidad & Son lecturas y escrituras sobre la cuenta. Cualquier arquitectura que cumpla ASR-08 los resuelve sin cambiar ningún componente. \\ \hline
Tienda, cajas y monedas (CU-38 a CU-40) & Funcionalidad & Necesitan transacciones, que SQL Server ya ofrece y que ASR-05 ya obliga a tratar como escritura confirmada. No piden ninguna decisión estructural nueva. \\ \hline
Amigos e invitaciones (CU-30 a CU-32, CU-49) & Funcionalidad & Usan el mismo servidor, comunicación y persistencia que ASR-08; el canal de amigos del chat ya entra en ASR-11. \\ \hline
Ranking e historial (CU-35, CU-36) & Funcionalidad & Son consultas sobre datos que otros ASR ya obligan a guardar. \\ \hline
Repetición de una partida (CU-37) & Funcionalidad & Se apoya en el determinismo de ASR-06 y en las jugadas registradas de ASR-05; no añade una decisión propia. \\ \hline
Tutorial (CU-41) & Funcionalidad & Es importante para REQ-07, pero se resuelve dentro del cliente y del motor que otros ASR ya exigen. \\ \hline
Moderadores y bitácoras (CU-46 a CU-48) & Funcionalidad & Son operaciones administrativas sobre datos que ASR-02 y ASR-11 ya obligan a registrar. \\ \hline
Originalidad frente a Quoridor y licencias de recursos (CON-15, CON-18; QAS-15) & Restricción & Importancia alta y sensibilidad baja: se cumplen con marca, textos y recursos fuera del código, una convención que ASR-13 ya impone. Ninguna estructura razonable lo impide. \\ \hline
Instalación en plataformas de distribución (Q-10, QAS-18, DEP-08) & Restricción condicionada & Solo aplica si se publica, y se resuelve con el empaquetado, sin tocar componentes. \\ \hline
DEC-01: modelo cliente-servidor basado en servicios & Decisión del equipo & No es un requisito, sino la respuesta del equipo a ASR-08 y ASR-02. Se registra aquí como decisión ya tomada. \\ \hline
DEC-02: no se confía en el cliente; toda la lógica del juego se ejecuta en el servidor & Decisión del equipo & Responde a ASR-02 de la forma más estricta posible. Condiciona cómo se cumplen ASR-01 y ASR-06, pero no es un requisito. \\ \hline
DEC-03: una partida en línea exige conexión con la base de datos & Decisión del equipo & Fija cómo se cumple ASR-05: ante una caída, la partida no continúa sin base de datos. Es una decisión sobre el requisito, no un requisito nuevo. \\ \hline
\end{tabla}

% ------------------------------------------------------------
\subsection{Priorización}\label{sec:asr:prioridad}

Ya confirmado el conjunto, falta decidir qué ASR pesa más. La priorización sigue el mismo esquema que el Utility Tree, con dos criterios independientes que se califican por separado y se combinan en un nivel de prioridad.

\textbf{Cómo se prioriza.} La \emph{importancia} mide cuánto le importa el requisito a quien puede decidir sobre el proyecto; la \emph{sensibilidad arquitectónica} mide cuánto cambia la estructura según se cumpla o no. La escala de cada criterio es la siguiente.

\begin{tabla}{|L{2.8cm}|Y|Y|}
\hline
\filacab \cab{Nivel} & \cab{Importancia} & \cab{Sensibilidad arquitectónica} \\ \hline
\endhead
Alta (H) & Lo impone el cliente, la regulación o un requisito impuesto, o su incumplimiento provoca el rechazo de un stakeholder con poder de decisión. & Hay alternativas estructurales razonables y solo algunas lo cumplen, y cambiar de alternativa después afecta a varios componentes a la vez. \\ \hline
Media (M) & Depende de una pregunta abierta o de una funcionalidad valiosa pero no esencial. & Su efecto queda dentro de un componente, o en convenciones que atraviesan el sistema sin cambiar su partición. \\ \hline
Baja (L) & Es deseable, pero ningún stakeholder lo exige. & Se cumple igual en cualquier alternativa razonable. \\ \hline
\end{tabla}

Las dos calificaciones se combinan en cuatro niveles. Un ASR (H, H) es \textbf{prioritario}: es un criterio de descarte, y una alternativa que no pueda cumplirlo se rechaza aunque sea mejor en todo lo demás. Un ASR (H, M) tiene prioridad \textbf{alta}: es obligatorio, pero se cumple dentro de la estructura que los prioritarios imponen. Un ASR (M, M) tiene prioridad \textbf{media}. Un ASR con importancia media y sensibilidad alta queda \textbf{reservado} mientras dependa de una pregunta abierta, porque tomarlo en cuenta hoy obligaría a diseñar para una situación que puede no llegar. Ningún ASR resultó con calificación baja: un requisito así no habría pasado la prueba. Dentro de cada nivel, el orden sigue a las dependencias: primero va lo que los demás dan por hecho.

\textbf{ASR priorizados.} La tabla ordena los quince ASR con la justificación de cada una de sus dos calificaciones.

\begin{tabla}{|L{1.05cm}|L{1.3cm}|Y|Y|L{1.9cm}|}
\hline
\filacab \cab{Orden} & \cab{ASR} & \cab{Importancia y por qué} & \cab{Sensibilidad y por qué} & \cab{Prioridad} \\ \hline
\endhead
1 & ASR-08 & H. Lo impone el cliente (REQ-01) y es el objetivo central del juego. & H. Sin él no hay servidor, red ni protocolo: es la diferencia estructural más grande de la lista. & Prioritario \\ \hline
2 & ASR-02 & H. Lo exige STK-15, que da por hecho un cliente alterado (CRN-23). & H. Separa el servidor que decide del que solo retransmite, y ese cambio rehace cliente, servidor y protocolo. & Prioritario \\ \hline
3 & ASR-14 & H. La tecnología la asigna el cliente (STK-03) y no se negocia. & H. Obliga a un protocolo propio y a un modelo relacional, que otras tecnologías darían resueltos. & Prioritario \\ \hline
4 & ASR-03 & H. Sin él, el segundo factor de REQ-03 no protege nada y el chat de menores viaja legible (STK-15, STK-02). & H. Dónde vive el cifrado cambia el formato de los mensajes en los dos extremos. & Prioritario \\ \hline
5 & ASR-04 & H. Los jugadores son niños en redes domésticas inestables (CRN-06, STK-01). & H. Obliga a separar la sesión de partida de la conexión, lo que toca servidor, protocolo y cliente. & Prioritario \\ \hline
6 & ASR-01 & H. Es el único requisito con un número impuesto (REQ-02). & H. Su reparto cambia el flujo de la jugada en todos los componentes que atraviesa. & Prioritario \\ \hline
7 & ASR-05 & H. Lo exige STK-13 (CRN-19), sobre una base de datos alojada por un tercero (DEP-03, DEP-04). & H. Decide qué componente tiene el estado vigente y cuándo se escribe. & Prioritario \\ \hline
8 & ASR-06 & H. El diseñador cambiará las reglas decenas de veces (CRN-13) y CON-15 puede obligar a variantes. & H. Decide si el motor es una pieza independiente y determinista o queda mezclado con el servidor y la interfaz. & Prioritario \\ \hline
9 & ASR-11 & H. Lo imponen REQ-04 y REQ-05, y STK-02 puede vetar el juego. & M. Su efecto queda en el chat y la moderación, una vez decidido quién tiene la autoridad. & Alta \\ \hline
10 & ASR-12 & H. Lo imponen REQ-03 y REQ-07. & M. Queda dentro del módulo de cuentas y de sus dependencias externas. & Alta \\ \hline
11 & ASR-13 & H. Lo imponen REQ-06 y el objetivo de distribución internacional. & M. Impone convenciones transversales, no una partición distinta. & Alta \\ \hline
12 & ASR-15 & H. El plazo y los incrementos los fijan el cliente y el auditor (STK-03, STK-16). & M. Acota las alternativas sin elegir entre las razonables. & Alta \\ \hline
13 & ASR-09 & M. Es funcionalidad valiosa pero no esencial. & M. Pide un mecanismo de difusión dentro del servidor. & Media \\ \hline
14 & ASR-10 & M. CU-27 lo especifica, pero Q-01 sigue abierta en el documento de contexto. & M. La IA entra como un participante más dentro de la estructura de partidas. & Media \\ \hline
15 & ASR-07 & M. Depende de Q-02. & H. Si se aprueba, obliga a separar el servidor y su configuración. & Reservado \\ \hline
\end{tabla}

\textbf{Por qué este orden entre los prioritarios.} ASR-08 va primero porque sin él no hay servidor que arbitrar; ASR-02 decide que ese servidor tiene la autoridad; ASR-14 fija con qué se construye todo. Los cinco restantes piden al servidor, a la comunicación y a la persistencia comportamientos que solo tienen sentido una vez que los tres primeros existen: la protección del tráfico y la supervivencia a una desconexión se apoyan en el canal que ASR-14 obliga a construir, y el tiempo de respuesta, la persistencia y las reglas, en el servidor que ASR-08 y ASR-02 ya establecieron.

% ------------------------------------------------------------
\subsection{Resultado}

Quedan catorce ASR activos y uno reservado. Los ocho prioritarios, ASR-01 a ASR-06, ASR-08 y ASR-14, son el criterio de descarte de cualquier alternativa de diseño. Los cuatro de prioridad alta, ASR-11 a ASR-13 y ASR-15, son obligatorios pero se cumplen dentro de la estructura que los prioritarios imponen. Los dos de prioridad media, ASR-09 y ASR-10, dependen de funcionalidades no esenciales o de una pregunta abierta. Junto con las tres decisiones del equipo, este conjunto es la entrada de todo lo que sigue en el documento.
